\section{Path-order curvature}

Let $A,B\in\operatorname{End}(\X)$ and let $s,t\in\R\setminus\{0\}$. Define two local actions
\begin{equation}
U_s=I+sA,\qquad V_t=I+tB.
\label{eq:actions}
\end{equation}
For $x\in\X$, compare the two orderings
\begin{equation}
x_{A\to B}=V_tU_sx,\qquad
x_{B\to A}=U_sV_tx.
\label{eq:orders}
\end{equation}

\begin{theorem}[T6: exact path-order curvature]
The ordered-path residue is
\begin{equation}
x_{A\to B}-x_{B\to A}=st(BA-AB)x.
\label{eq:commutator-residue}
\end{equation}
Thus the commutator
\begin{equation}
\mathcal K_{A,B}=BA-AB
\label{eq:curvature}
\end{equation}
is the exact first path-order curvature of the two local actions.
\end{theorem}

\begin{proof}
Expand both compositions:
\[
V_tU_sx=x+sAx+tBx+stBAx,
\]
\[
U_sV_tx=x+tBx+sAx+stABx.
\]
Subtracting gives \cref{eq:commutator-residue}.
\end{proof}

\begin{corollary}[Path independence]
The two orderings agree on $x$ if and only if $\mathcal K_{A,B}x=0$.
\end{corollary}

\begin{proposition}[Unitary covariance and magnitude invariance]
Let $Q:\X\to\X$ be unitary and put
\[
A'=QAQ^*,\qquad B'=QBQ^*,\qquad x'=Qx.
\]
Then
\begin{equation}
\mathcal K_{A',B'}=Q\mathcal K_{A,B}Q^*,
\qquad
\norm{\mathcal K_{A',B'}x'}=\norm{\mathcal K_{A,B}x}.
\label{eq:curvature-unitary}
\end{equation}
\end{proposition}

\begin{proof}
Direct multiplication gives
\[
B'A'-A'B'
=QBQ^*QAQ^*-QAQ^*QBQ^*
=Q(BA-AB)Q^*.
\]
Applying this identity to $x'=Qx$ and using unitarity of $Q$ proves the norm equality.
\end{proof}

\begin{proposition}[Commutator consistency]
The path-order curvature is antisymmetric and satisfies the Jacobi identity:
\[
\mathcal K_{A,B}=-\mathcal K_{B,A},
\qquad
[A,[B,C]]+[B,[C,A]]+[C,[A,B]]=0.
\]
\end{proposition}

\begin{proof}
Both identities follow by expanding the commutators and cancelling equal products with opposite signs.
\end{proof}

\begin{remark}
The term ``path-order curvature'' refers only to the exact order defect of the local actions in \cref{eq:actions}. No identification with Riemannian curvature or the curvature of a particular connection is asserted without additional structure.
\end{remark}

\begin{definition}[Curvature magnitude]
Define
\begin{equation}
\I_{\mathrm{curv}}(x;A,B)=\norm{\mathcal K_{A,B}x}
=\frac{\norm{x_{A\to B}-x_{B\to A}}}{|st|}.
\label{eq:curvature-magnitude}
\end{equation}
\end{definition}

\begin{proposition}[Curvature-sensitive admissibility]
If a policy fixes $\I_{\mathrm{curv}}\le\varepsilon_{\mathrm{curv}}$, then any candidate with $\norm{\mathcal K_{A,B}x}>\varepsilon_{\mathrm{curv}}$ cannot be identified with a path-independent accepted state under that policy.
\end{proposition}

\begin{proof}
Such an identification would contradict the declared inequality in \cref{eq:curvature-magnitude}.
\end{proof}

\begin{example}[Exact nonzero residue]
Let
\[
A=\begin{pmatrix}0&1\\0&0\end{pmatrix},\quad
B=\begin{pmatrix}0&0\\1&0\end{pmatrix},\quad
x=\begin{pmatrix}1\\1\end{pmatrix},\quad
s=\frac12,\quad t=\frac13.
\]
Then
\[
(BA-AB)x=\begin{pmatrix}-1\\1\end{pmatrix},
\qquad
x_{A\to B}-x_{B\to A}=\begin{pmatrix}-1/6\\1/6\end{pmatrix}.
\]
If $B=\lambda A$, the commutator vanishes and the two orderings agree exactly.
\end{example}

\section{Finite proof admission}

We now formalize a proof-shaped event without reference to prose quality.

\begin{definition}[Finite derivation object]
A finite derivation object is
\begin{equation}
\Pi=(\Gamma,\varphi_{\mathrm{target}},(\ell_i)_{i=1}^n,D,\rho),
\label{eq:proof-object}
\end{equation}
where $\Gamma$ is a finite set of premises, $\varphi_{\mathrm{target}}$ the declared target, $D$ a directed dependency graph on the proof lines, and $\rho$ a fixed finite collection of admissible inference rules.
\end{definition}

\begin{definition}[Admitted derivation]
The object $\Pi$ is admitted when:
\begin{enumerate}[label=(\roman*)]
\item $D$ is acyclic;
\item every dependency of line $\ell_i$ points to a predecessor of $\ell_i$;
\item every line is either a declared premise or follows from its cited predecessors by one rule in $\rho$;
\item the final verified line equals $\varphi_{\mathrm{target}}$.
\end{enumerate}
\end{definition}

\begin{theorem}[T7: finite proof admission]
If $\Pi$ is admitted, then $\Gamma\vdash_{\rho}\varphi_{\mathrm{target}}$. Conversely, any derivation with a cycle, an unlicensed rule, a missing dependency, or a final-target mismatch fails the admission predicate.
\end{theorem}

\begin{proof}
Because $D$ is acyclic, choose a topological ordering of its lines. Induct along that ordering. A premise line belongs to $\Gamma$. Every non-premise line follows from already verified predecessors by a rule in $\rho$, so each line is derivable from $\Gamma$. The final verified line is exactly $\varphi_{\mathrm{target}}$, proving $\Gamma\vdash_{\rho}\varphi_{\mathrm{target}}$. The converse statements follow directly because each listed defect violates one defining clause of admission.
\end{proof}

\begin{theorem}[T7B: semantic soundness lift]
Let $\mathfrak M$ be a semantics for the formulas of $\rho$. Assume every premise in $\Gamma$ is true in $\mathfrak M$ and every rule in $\rho$ is truth-preserving in $\mathfrak M$. If $\Pi$ is admitted, then
\[
\mathfrak M\models\varphi_{\mathrm{target}}.
\]
\end{theorem}

\begin{proof}
Use the same topological ordering as in T7. Every premise line is true by assumption. If all cited predecessors of a non-premise line are true, truth preservation of the licensed rule makes that line true. Induction therefore makes every verified line true, including the final line $\varphi_{\mathrm{target}}$.
\end{proof}

\begin{remark}
T7 is syntactic and relative to the declared calculus $\rho$. T7B is the semantic lift: it requires sound rules and true premises. Neither theorem claims that $\rho$ is complete for all mathematics.
\end{remark}
